
%%%%%%%%%%%%%%%%%%%%%%%%%%%%%%%%%%%%%%%%
%           main body
%%%%%%%%%%%%%%%%%%%%%%%%%%%%%%%%%%%%%%%%


%-----------------------------------
%	TITLE PAGE
%-----------------------------------

\title[Linear Algebra]{\LARGE \bfseries 第二部分\quad 线性代数} % The short title appears at the bottom of every slide, the full title is only on the title page
\author[Eigenvalue and Eigenvector] % Your institution as it will appear on the bottom of every slide, may be shorthand to save space
{\Large \bfseries \S 5. 特征值与特征向量}
% \author{Singular Value Decomposition} % Your name
% \date{\today} % Date, can be changed to a custom date
\date{}

%------------------------------------------------

\begin{document}


%------------------------------------------------

\begin{frame}

\titlepage % Print the title page as the first slide

\end{frame}

%------------------------------------------------

\begin{frame}
\frametitle{内容提要} % Table of contents slide, comment this block out to remove it
\tableofcontents % Throughout your presentation, if you choose to use \section{} and \subsection{} commands, these will automatically be printed on this slide as an overview of your presentation
\end{frame}

%-----------------------------------
%	PRESENTATION SLIDES
%-----------------------------------

%------------------------------------------------

\section{特征值与特征向量}

%------------------------------------------------

\begin{frame}

\begin{block}{问题引入}
之前已经讲过，一个线性变换$\varphi: V \to V$对应于一个方阵
\[\varphi(e_1, \ldots, e_n) = (e_1, \ldots, e_n) \underbrace{\begin{pmatrix}
a_{11} & \cdots & a_{1n} \\ \vdots & & \vdots \\ a_{n1} & \cdots & a_{nn} \end{pmatrix}}_{A}\]

\pause

但是这样的表示形式稍显复杂，希望得到$\varphi$的尽可能简单的表示。

\pause

可能的办法：换一组基，$\varepsilon_1, \ldots, \varepsilon_n$, 使得
\[\varphi(\varepsilon_1) = \lambda_1 \varepsilon_1, \; \ldots, \; \varphi(\varepsilon_n) = \lambda_n\varepsilon_n\]
也就是说，此时$\varphi$对应的方阵为$diag(\lambda_1 , \ldots, \lambda_n)$.
\end{block}

\end{frame}

%------------------------------------------------

\begin{frame}

\begin{block}{问题引入}
\vspace{1em}
于是，我们有如下的问题
\vspace{2em}
\begin{itemize}
\item<2-> 这样的$\varepsilon_1,\ldots,\varepsilon_n$是否存在?
\vspace{1em}
\item<3-> 如果存在，如何找这些$\varepsilon_i$以及$\lambda_i$?
\end{itemize}
\end{block}

\end{frame}

%------------------------------------------------

\begin{frame}

\begin{block}{特征值与特征向量}
如果有$\lambda \in \mathbb{R}$以及{\color{red}非零向量}$\alpha \in V$, 使得
\[\varphi(\alpha) = \lambda \alpha,\]
那么称$\lambda$为线性变换$\varphi$的特征值（Eigenvalue或Characteristic Value），向量$\alpha$被称为线性变换$\varphi$的属于特征值$\lambda$的特征向量（Eigenvector）。
\end{block}

\end{frame}

%------------------------------------------------

\begin{frame}

\begin{block}{特征值与特征向量}
假设我们从$V$的某一组基$e_1, \ldots, e_n$出发，开始寻找$\varphi$的特征值与特征向量。

\vspace{1em}
\pause

设$\varphi$在这组基下的矩阵表达为$A$. 对于向量$\alpha \in V$, 设它在这组基下的坐标为$\mathbf{x} = (x_1,\ldots,x_n)^T$.于是，如果$\alpha$是线性变换$\varphi$的属于特征值$\lambda$的特征向量，那么
\[A\mathbf{x} = \lambda\mathbf{x}\]
于是，$\lambda$也被称为方阵$A$的特征值，列向量$\mathbf{x}$被称为方阵$A$的的属于特征值$\lambda$的特征列向量。


\vspace{1em}
\pause
\begin{itemize}
\item 特征值与基的选取无关；相似的方阵有同一组特征值。
\item 特征向量有很多：设$\alpha_1,\alpha_2$都是属于$\lambda$的特征向量，那么任取不同时为零的实数$t_1,t_2$, 新的向量$t_1\alpha_1+t_2\alpha_2$也是属于$\lambda$的特征向量。
\end{itemize}
\end{block}

\end{frame}

%------------------------------------------------

\begin{frame}

\begin{block}{如何求特征值}
设一个实数$\lambda\in\mathbb{R}$是方阵$A$的特征值，即存在非零列向量$\mathbf{x}$使得$A\mathbf{x} = \lambda\mathbf{x}$. 那么
\begin{align*}
A\mathbf{x} = \lambda\mathbf{x} & \Longleftrightarrow (A - \lambda I_n) \mathbf{x} = \mathbf{0} \\
& \Longrightarrow \det (A - \lambda I_n) = 0
\end{align*}

\pause

反过来，如果$\det (A - \lambda I_n) = 0$，那么线性方程组$(A - \lambda I_n) \mathbf{x} = \mathbf{0}$总是有（无穷多）解的。任取一个非零解$\mathbf{x}_0$, 它都满足$(A - \lambda I_n) \mathbf{x}_0 = \mathbf{0}$, 即有$A\mathbf{x}_0 = \lambda\mathbf{x}_0$. 于是我们知道$\lambda$是方阵$A$的特征值，而$\mathbf{x}_0$是属于$\lambda$的特征向量。
\end{block}

\end{frame}

%------------------------------------------------

\begin{frame}

\begin{block}{特征方程}
称以$\lambda$为未定元的多项式
\[f(\lambda) := \det(\lambda I_n - A) \in \mathbb{R}[\lambda]\]
为方阵$A$的特征多项式（Characteristic Polynomial）。

\vspace{2em}
\pause

\begin{itemize}
\item 方阵$A$的特征值，就是方阵$A$的特征多项式$f(\lambda)$的根。
\end{itemize}
\end{block}

\end{frame}

%------------------------------------------------

\begin{frame}

\begin{block}{特征方程}
\begin{itemize}
\item 回忆一下，我们之前讲过，只要一个数学实体$R$有``加法''、``乘法''以及``数乘''，我们就能定义元素为$R$中元素的$m\times n$矩阵，其全体记作$M_{m\times n}(R)$. 其方阵$(a_{ij})_{n\times n}$的行列式可以依照实数方阵的行列式类似定义：
\[\det(a_{ij}) = \sum\limits_{j_1\ldots j_n} (-1)^{\tau(j_1\ldots j_n)} a_{1j_1}\cdots a_{nj_n}\]
\item 对于任意的一个可逆方阵$P$，有
\begin{align*}
\det(\lambda I_n - P^{-1}AP) = \det (P^{-1}(\lambda I_n - A)P) = \det(\lambda I_n - A)
\end{align*}
也就是说相似方阵有相同的特征多项式。对于线性变换$\varphi$来说，无论$V$取哪一组集，得到相应的方阵表示，它们的特征多项式都是一样的。于是，$f(\lambda)$也被称为线性变换$\varphi$的特征多项式。
\end{itemize}
\end{block}

\end{frame}

%------------------------------------------------

\begin{frame}

\begin{block}{没有特征值的情况}
{\bfseries 注意}！一个实方阵不一定有实的特征值。

\vspace{1em}
\pause

例如，考虑2阶方阵$A = \begin{pmatrix} \cos\theta & -\sin\theta \\ \sin\theta & \cos\theta \end{pmatrix}$, 它的特征多项式为
\begin{align*}
f(\lambda) &= \det (\lambda I_2 - A) = \det \begin{pmatrix} \lambda - \cos\theta & \sin\theta \\ -\sin\theta & \lambda - \cos\theta \end{pmatrix} \\
& = (\lambda - \cos\theta)(\lambda - \cos\theta) + \sin\theta\sin\theta \\
& = \lambda^2 - 2\cos\theta\lambda + 1
\end{align*}
以上方程当$\cos\theta\neq 1$时是没有实根的。

\vspace{1em}
\pause

方阵$A$对应于2维实平面$\mathbb{R}^2$中的逆时针旋转$\theta$角的线性变换。除非这个角度是$2\pi$的倍数，否则不可能有向量经过这个变换之后还保持方向不变。
\end{block}

\end{frame}

%------------------------------------------------

\begin{frame}

\begin{block}{如何求特征向量}
从方阵$A$的特征多项式$f(\lambda)$出发，算得其特征值$\lambda_1, \ldots, \lambda_s$, 接下来就可以从$s$个齐次线性方程组出发
\[(A - \lambda_i I_n)\mathbf{x} = \mathbf{0}\]
算得属于相应特征值的特征向量。
\end{block}

\end{frame}

%------------------------------------------------

\begin{frame}

\begin{block}{几何重数与代数重数}
对于方阵$A$的一个特征值$\lambda_i$, 称
\[V_{\lambda_i} := \{ \mathbf{x} \in \mathbb{R}^n \ |\ A\mathbf{x} = \lambda_i\mathbf{x} \}\]
为$A$的关于特征值$\lambda_i$的特征子空间，其维数为特征值$\lambda_i$的几何重数。注意，零向量满足上式，自然包含在所有的$V_{\lambda_i}$中，但特征向量不能是零向量。也就是说
\[V_{\lambda_i} := \{ \text{属于$\lambda_i$的特征向量} \} \cup \{\mathbf{0}\}\]

\vspace{1em}
\pause

相应地，特征多项式$f(\lambda)$的因子$(\lambda - \lambda_i)$的次数被称为特征值$\lambda_i$的代数重数。
\end{block}

\end{frame}

%------------------------------------------------

\begin{frame}

\begin{block}{算例}
求方阵$A = \begin{pmatrix} 2 & 2 & -1 \\ -1 & 0 & 1 \\ -1 & 0 & 2 \end{pmatrix}$的特征值与特征多项式。

\vspace{1em}
\pause

解：先求$A$的特征值，特征多项式为
\[
f(\lambda) = \det \begin{pmatrix} \lambda-2 & -2 & 1 \\ 1 & \lambda & -1 \\ 1 & 0 & \lambda-2 \end{pmatrix} = (\lambda-1)^2(\lambda-2)
\]
所以得到$A$的特征值为：$\lambda_1=1$（重数为2）；$\lambda_2=2$。
\end{block}

\end{frame}

%------------------------------------------------

\begin{frame}

\begin{block}{算例}
再求$A$的特征向量。

\begin{itemize}
\item 对于$\lambda_1 = 1$，解齐次线性方程组
\[(\lambda_1 I-A)\mathbf{x} = \mathbf{0},\]
得到属于特征值$\lambda_1 = 1$的一个线性无关的特征向量为$\mathbf{x}_{11} = t(1,0,1)^T, t\in\mathbb{R}, t\neq 0$.
\item 对于$\lambda_2 = 2$，解齐次线性方程组
\[(\lambda_2 I-A)\mathbf{x} = \mathbf{0},\]
得到属于特征值$\lambda_2 = 2$的一个线性无关的特征向量为$\mathbf{x}_{21} = t(0,1,2)^T, t\in\mathbb{R}, t\neq 0$.
\end{itemize}
\end{block}

\end{frame}

%------------------------------------------------

\section{特征值与特征向量的性质}

%------------------------------------------------

\begin{frame}

\begin{block}{特征值与特征向量的性质}
\begin{itemize}
\item<1-> $n$阶方阵的特征值至多有$n$个，可以少于$n$个。
\item<2-> 每个特征值都对应有至少一个特征向量，但其特征空间的维数（几何重数）不超过其代数重数，可以严格小于。
\item<3-> 设$f(\lambda) = \det(\lambda I_n - A) = \lambda^n - a_1\lambda^{n-1} + \cdots + (-1)^na_n$是方阵$A$的特征多项式，那么
\[a_1 = tr A, \quad a_n = \det A\]
\item<4-> 对于分块方阵\(\displaystyle A = \begin{pmatrix} A_1 & A_3 \\ 0 & A_2 \end{pmatrix}\), 其中$A_1,A_2$都是方阵，那么
\[f_A(\lambda) = f_{A_1}(\lambda)f_{A_2}(\lambda)\]
\end{itemize}
\end{block}

\end{frame}

%------------------------------------------------

\section{方阵的对角化}

%------------------------------------------------

\begin{frame}

\begin{block}{方阵的对角化}
\begin{itemize}
\item 方阵$A$相似于上三角阵\(\displaystyle \begin{pmatrix} \lambda_1 & & * \\ & \ddots & \\ & & \lambda_n \end{pmatrix}\)当且仅当其特征多项式$f(\lambda)$有$n$个特征值（可以重复）$\lambda_1,\ldots,\lambda_n$.

\vspace{2em}
\pause

\item 方阵可对角化的充要条件：\newline
方阵$A$相似于对角阵\(\displaystyle \begin{pmatrix} \lambda_1 & & \\ & \ddots & \\ & & \lambda_n \end{pmatrix}\)当且仅当他有$n$个线性无关的特征向量。此时，$\lambda_1,\ldots,\lambda_n$为$A$的特征值。

\vspace{1em}
\pause

因此，可对角化方阵的对角化又被称作方阵的特征分解（Eigendecomposition）。
\end{itemize}
\end{block}

\end{frame}

%------------------------------------------------

\begin{frame}

\begin{block}{方阵的对角化}
{\bfseries 引理}：设$\lambda_1\neq\lambda_2$是方阵$A$的两个不相等的特征值。任取对应于$\lambda_1$的一个特征向量$\alpha_1$, 任取对应于$\lambda_2$的一个特征向量$\alpha_2$, 那么$\alpha_1$与$\alpha_2$线性无关。

\vspace{2em}
\pause

{\bfseries 推论}：若$n$阶方阵有$n$个互不相等的特征值，那么该方阵可以对角化。
\end{block}

\end{frame}

%------------------------------------------------

\begin{frame}

\begin{block}{引理的证明}
有定义，$\alpha_1$与$\alpha_2$都不是零向量，因此假如他们线性相关，就可以设存在实数$k$使得$\alpha_1=k\alpha_2$. 那么
\[A\alpha_1 = A(k\alpha_2) = kA\alpha_2 = k\lambda_2\alpha_2.\]
另一方面，又有
\[A\alpha_1 = \lambda_1\alpha_1 = \lambda_1k\alpha_2.\]
于是
\[k\lambda_2\alpha_2 = k\lambda_1\alpha_2,\]
即
\[k(\lambda_1-\lambda_2)\alpha_2 = \mathbf{0}.\]
这与特征向量的定义（都是非零向量）矛盾。
\end{block}

\end{frame}

%------------------------------------------------

\begin{frame}

\begin{block}{方阵对角化的过渡矩阵}
设$n$阶方阵$A$有$n$个线性无关的特征列向量$\alpha_1,\ldots,\alpha_n$, 分别对应于特征值$\lambda_1,\ldots,\lambda_n$, 那么有
\[A(\alpha_1,\ldots,\alpha_n) = (\alpha_1,\ldots,\alpha_n) \underbrace{\begin{pmatrix} \lambda_1 & & \\ & \ddots & \\ & & \lambda_n \end{pmatrix}}_{\Lambda}.\]

\vspace{1em}
\pause

令$P = (\alpha_1,\ldots,\alpha_n)$, 为一个可逆$n$阶方阵，而且有$AP = P\Lambda,$ 即有
\[A = P\Lambda P^{-1}, \quad \text{ 或者说 } P^{-1}AP = \Lambda.\]
\end{block}

\end{frame}

%------------------------------------------------

\begin{frame}

\begin{block}{例题}
令$A = \begin{pmatrix} 2 & -1 & 2 \\ 0 & 3 & 1 \\ 0 & 0 & 5 \end{pmatrix}$, 求$A^{n}$.

\vspace{1em}
\pause

解：容易算得$A$的特征值为$2, 3, 5$, 对应的特征向量为\[k(1,0,0)^T, \; k(-1,1,0)^T, \; k(1,1,2)^T\]
那么令$P=\begin{pmatrix} 1 & -1 & 1 \\ 0 & 1 & 1 \\ 0 & 0 & 2 \end{pmatrix}$，则有$A = P \begin{pmatrix} 2 & 0 & 0 \\ 0 & 3 & 0 \\ 0 & 0 & 5 \end{pmatrix} P^{-1}$.
因此
\begin{equation*}
\small
A^{n} = P \begin{pmatrix} 2^{n} & 0 & 0 \\ 0 & 3^{n} & 0 \\ 0 & 0 & 5^{n} \end{pmatrix} P^{-1} = \begin{pmatrix} 2^{n} & 2^n-3^n & \frac{1}{2}(3^n+5^n-2^{n+1}) \\ 0 & 3^n & \frac{1}{2}(5^n-3^n) \\ 0 & 0 & 5^n \end{pmatrix}
\end{equation*}
\end{block}

\end{frame}

%------------------------------------------------

% \begin{frame}

% \begin{block}{Vandermonde行列式的计算}
% \begin{enumerate}
% % \addtocounter{enumi}{0}
% \item 把第$n-1$行乘以$-a_1$加到第$n$行，再把第$n-2$行乘以$-a_1$加到第$n-1$行。按照这种顺序依次向上，直到用第$2$行乘以$-a_1$加到第$1$行。由行列式性质，$V_n$的值不变。
% \[
% V_n = \det \begin{pmatrix}
% 1 & 1 & \cdots & 1 \\
% 0 & a_2-a_1 & \cdots & a_n-a_1 \\
% \vdots & \vdots & & \vdots \\
% 0 & a_2^{n-2}(a_2-a_1) & \cdots & a_n^{n-2}(a_n-a_1) \\
% \end{pmatrix}
% \]
% \end{enumerate}
% \end{block}

% \end{frame}

%------------------------------------------------
% % closing page
% \begin{frame}
% \HUGE{\centerline{\bfseries {\color{gray} The} {\color{zkblue} End}}}

% \pause

% \vspace{1.2em}

% \Large{\centerline{\bfseries {\color{gray}Thank} \quad {\color{zkblue} You}}}

% \end{frame}

%----------------------------------------------------------------------------------------

\end{document}
